%!TeX program = xelatex
\documentclass[fleqn]{article}
\usepackage[utf8]{inputenc}
\usepackage{amsmath}
\usepackage{graphicx}
\usepackage{mathtools}
\usepackage{datetime}
%\usepackage{fancyhdr}
\usepackage[english]{babel}
\usepackage{titlesec}
\usepackage{xcolor}
\usepackage[object=vectorian]{pgfornament}
\usepackage{lastpage}
\usepackage{fontspec}
\usepackage{tikzpagenodes}
\usepackage{stackengine}
\setstackEOL{\\}
\usepackage{geometry}
\setmainfont{GFSArtemisia.otf}

% Settings
\newcommand\class{DIT}
\newcommand\classSubject{Practica Basis Electronica}
\newcommand\user{Vincent van Setten}
\newcommand\week{Week 1}
\newcommand\lecturer{Marius Versteegen}

% Theming
\oddsidemargin = 2pt
\geometry{
    top=3in,
    inner=0.5in,
    outer=0.5in,
    bottom=1in,
    headheight=7ex,
    headsep=6ex
}
\definecolor{hublue}{RGB}{35, 136, 196}
\definecolor{hured}{RGB}{214, 30, 37}

\titleformat*{\section}{\Large\bfseries\fontfamily{cmr}\selectfont\color{hublue} }

%% Don't Touch  %%
\newdateformat{monthyeardate}{%
  \monthname[\THEMONTH], \THEYEAR}

\newpagestyle{hubanner}{

\headrule
\sethead{}{
\Huge\fontfamily{cmr}\selectfont\color{white}\sectiontitle 
    }{}
}
\pagestyle{hubanner}


\begin{document} 

\renewcommand\makeheadrule{%
\begin{tikzpicture}[remember picture,overlay]
  \fill[hublue] 
  (-4,0) 
    rectangle
  (25, 20);
\end{tikzpicture}}

\begin{titlepage} 

{\fontfamily{cmr}\selectfont
	\newcommand{\HRule}{\rule{\linewidth}{0.5mm}}
	
	\center
	
	%------------------------------------------------
	%	Headings
	%------------------------------------------------
	
	\textsc{\LARGE HU University of Applied Sciences Utrecht}\\[1.5cm] 
	
	\textsc{\Large \class}\\[0.5cm] 
	
	\textsc{\large \week}\\[0.5cm]
	
	%------------------------------------------------
	%	Title
	%------------------------------------------------
	
	\HRule\\[0.4cm]
	
	{\huge\bfseries \classSubject}\\[0.4cm] % Title of your document
	
	\HRule\\[1.5cm]
	
	%------------------------------------------------
	%	Author(s)
	%------------------------------------------------
	
	\begin{minipage}{0.4\textwidth}
		\begin{flushleft}
			\large
			\textit{Student}\\
			\user % Your name
		\end{flushleft}
	\end{minipage}
	~
	\begin{minipage}{0.4\textwidth}
		\begin{flushright}
			\large
			\textit{Teacher}\\
			\lecturer % Supervisor's name
		\end{flushright}
	\end{minipage}
	
	%------------------------------------------------
	%	Date
	%------------------------------------------------
	
	\vfill\vfill\vfill % Position the date 3/4 down the remaining page
	
	{\large\monthyeardate\today}% Date, change the \today to a set date if you want to be precise
	
	%------------------------------------------------
	%	Logo
	%------------------------------------------------
	
	\vfill
	\vfill
	 \includegraphics[width=0.2\textwidth]{../images/hu.png}
	%----------------------------------------------------------------------------------------
	
     }
\end{titlepage}

%% DO TOUCH %%
\chapter{Introduction}
Het vak \class verteld de basis van elektrotechniek. Op deze manier worden we bekend met het gebruik van basis
elektronische componenten en hoe we daarmee berekeningen uitvoeren.
Hoewel ik thuis wel een aantal van deze componenten heb, heb ik niet genoeg spullen om veel van deze opdrachten fysiek
uit te voeren. Onder andere is mijn multimeter overleden(hij stond per ongeluk een erg lange tijd aan op de
detectiemodus). 
Daarnaast heb ik dit huiswerk al een keer eerder gedaan(inclusief de eindoprachten). Ik maak dit huiswerk dus vooral ter
voorbereiding op het tentamen. Ik zal proberen het huiswerk daarentegen wel zo goed mogelijk te maken. 




\end{document}
